%=======================================================================
\section{Results}
\label{sec:results}
%=======================================================================
% SKELETON. Subsections mirror the IHTC-18 paper. The numbers in
% Table~\ref{tab:ref-results} were produced with thums.py v0.13 at the case of
% Table~\ref{tab:conditions} (n_s = 100, which the code flags); regenerate
% them with the final numerics before use.

\subsection{Power-to-heat-to-power cycle analysis}
\label{sec:res-cycles}

\todo{$T$--$s$ diagram of both cycles at the reference case; table of state
points, pressures, $y$, $x_\mathrm{8h}$, $\varepsilon_\mathrm{R1}$,
$\varepsilon_\mathrm{R2}$, $\etaORC$, $\COP$, duties and flow rates.}
\todo{ORC evaporator composite curves showing the interior pinch at the onset
of boiling, with the hot-end rule overlaid: in the first pass $T_\mathrm{3e}$
falls from \SI{134.1}{\celsius} to \SI{95.1}{\celsius} and $\etaORC$ from 0.198
to 0.148 (v0.12, irreversible machines; the corresponding reversible values are
0.232 and 0.174).}
\todo{Cost of machine irreversibility: rerun with
$\eta_{s,t}=\eta_{s,p}=\eta_{s,c}=1$ (recovers $\etaORC=0.1738$, $\COP=2.759$)
and report the change in RTE; basis for a component-wise exergy balance.}
\todo{$UA$ of the four plant exchangers; the ORC condenser dominates.}

\subsection{Thermal analysis of the borehole regenerator at cyclic steady state}
\label{sec:res-BB}

\begin{table}[htbp]
\centering
\caption{Reference case at cyclic steady state (model v0.13, provisional).}
\label{tab:ref-results}
\small
\resizebox{\linewidth}{!}{%
\begin{tabular}{@{}l r l@{}}
\toprule
quantity & value & note \\
\midrule
$\COP$ / $\etaORC$ & 2.473 / 0.1487 & irreversible machines; corrected pass \\
$T_\mathrm{3e}$ & \SI{95.5}{\celsius} & pinch-limited (hot-end rule: \SI{136.5}{\celsius}) \\
$x_\mathrm{8h}$ / $y$ & 0.436 / 0.805 & \\
$\varepsilon_\mathrm{R1}$ / $\varepsilon_\mathrm{R2}$ & 0.205 / 0.120 & required, Eq.~\eqref{eq:eps-R} \\
HTHP suction superheat $T_\mathrm{1h}-T_\mathrm{13h}$ & \SI{14.3}{\kelvin} & \SI{26.5}{\kelvin} if isentropic \\
$\Nw$ & 15.62 & latent inventory, Eq.~\eqref{eq:Nwell} \\
$\md_\mathrm{ch}$ / $\md_\mathrm{dc}$ per hairpin & 0.965 / \SI{0.971}{\kilo\gram\per\second} & pinned by the glides, Eq.~\eqref{eq:mdot} \\
cycles to CSS & 73 & from a fully solid store \\
$\eta_\mathrm{storage}$ & 1.000000 & Eq.~\eqref{eq:css-unity} \\
deviation $\mathcal D$ & $+\SI{4.46}{\percent}$ & Eq.~\eqref{eq:deviation} \\
$\overline{T}_\mathrm{2d}$ / $\overline{T}_\mathrm{2c}$ & \SI{148.5}{\celsius} / \SI{102.5}{\celsius} & nominal 146.1 / 105.0 \\
$\bar\varepsilon_\mathrm{ch}$ / $\bar\varepsilon_\mathrm{dc}$ & 0.9704 / 0.0109 & cycled fraction 0.9595 \\
$\eta_T$ / $\eta_{T,\mathrm{eff}}$ & 0.1413 / 0.1365 & \\
RTE / $\varepsilon_\mathrm{RTE}$ & 0.3172 / 0.3044 & v0.11 (multipliers): 0.3003 \\
$\psi$ / $\psi_\mathrm{str}$ & 0.2356 / 0.2764 & Eq.~\eqref{eq:psi}, with pumping \\
$\md_\mathrm{geo}$ / $N_\mathrm{geo}$ & \SI{100.8}{\kilo\gram\per\second} / 8.07 & Eq.~\eqref{eq:mgeo}; $\md_\mathrm{geo,well}=\SI{12.5}{\kilo\gram\per\second}$ \\
$(\Nw+N_\mathrm{geo})/\dot W_\mathrm{el,out}$ & \SI{22.7}{\per\mega\watt} & storage plus geothermal wells \\
$\Delta E_\mathrm{therm}$ / $\Delta E_\mathrm{elec}$ per well & 4731 / \SI{646}{\kilo\watt\hour} & \\
$\rho_E$ & \SI{170.9}{\kilo\watt\hour\per\meter\cubed} & \\
pumping, charge / discharge & 34.9 / \SI{35.4}{\kilo\watt} & \SI{6.73}{\percent} of gross output \\
$UA$, four plant exchangers & \SI{2253}{\kilo\watt\per\kelvin} & ORC condenser 1135 \\
\bottomrule
\end{tabular}}
\end{table}

\chk{Grid sensitivity, v0.12 (not yet repeated on v0.13): with $n_s=99$ (layer boundaries on segment
boundaries) the deviation is $+\SI{4.44}{\percent}$ instead of
$+\SI{4.53}{\percent}$, $\bar\varepsilon_\mathrm{ch}$ is 0.980 instead of 0.972,
and the residual $\bar\varepsilon_\mathrm{dc}$ doubles, 0.025 instead of 0.012,
while RTE is unchanged to four figures. The residual melted fraction is the
diagnostic used to separate rate from inventory limitation, so the final runs
must use a layer-conforming grid and a convergence check.}

\todo{Profiles along the well at CSS (depth coordinate): $T$, $\Tpcm$ against
the $T_m$ staircase, $\eloc$ at several times in each half-cycle, $\Ui$ and
$\NTU$. Emphasise the sawtooth imposed by the cascade and where the PCM is
superheated or subcooled.}
\todo{Convergence to CSS: charged and discharged energy per cycle and state
drift versus cycle number.}
\todo{Decomposition of the energy per half-cycle into plateau and sensible
branches; peak superheat and subcooling.}
\todo{Resistance budget of $\Ui$ (convection, fouling, wall, PCM side) over a
half-cycle.}

\subsection{Exergy analysis}
\label{sec:res-exergy}

Table~\ref{tab:exergy} gives the exergy audit of the reference case over one
cycle. It is evaluated on the same basis as the round-trip efficiency: the
energy the field actually moved at cyclic steady state, with the borehole
parasitics charged as both an input and a destruction, so that $\psi$ and RTE
differ only in the charge made for the geothermal stream. Earlier drafts
evaluated it on the target duties and excluded the parasitics, which made the
two incomparable; the code was corrected to match this definition, not the
reverse. Three features stand out.

\begin{table}[htbp]
\centering
\caption{Exergy audit of the reference case over one cycle at CSS, resource
convention, $T_0=\SI{20}{\celsius}$ (model v0.13, provisional).}
\label{tab:exergy}
\small
\begin{tabular}{@{}l r r@{}}
\toprule
 & \si{\kilo\watt\hour} & share of $\sum I_j$ [\si{\percent}] \\
\midrule
\multicolumn{3}{@{}l}{\emph{inputs}}\\
electrical, $W_\mathrm{el,in}$ (incl.\ charge pumping) & 31813 & \\
geothermal, $\Ex_\mathrm{geo}$ (\SI{100.8}{\kilo\gram\per\second}) & 11019 & \\
\multicolumn{3}{@{}l}{\emph{outputs and loss}}\\
electrical, $W_\mathrm{el,out}$ & 10000 & \\
reinjected at \SI{50}{\celsius}, $\Ex_\mathrm{reinj}$ & 6326 & ($u_\mathrm{geo}=0.426$) \\
\midrule
\multicolumn{3}{@{}l}{\emph{destruction}}\\
HTHP condenser & 4584 & 17.4 \\
ORC evaporator & 3825 & 14.5 \\
HTHP throttle 7h$\to$8h & 3034 & 11.5 \\
HTHP compressor, high stage & 1915 & 7.3 \\
borehole field, Eq.~\eqref{eq:I-BB} & 1740 & 6.6 \\
HTHP motor & 1573 & 6.0 \\
ORC condenser & 1306 & 4.9 \\
HTHP evaporator & 1233 & 4.7 \\
HTHP compressor, low stage & 1222 & 4.6 \\
ORC turbine, low pressure & 1215 & 4.6 \\
HTHP throttle 9h$\to$4h & 1168 & 4.4 \\
ORC feed heater & 830 & 3.1 \\
borehole pumping & 703 & 2.7 \\
HTHP regenerator R2 & 671 & 2.5 \\
ORC generator & 550 & 2.1 \\
ORC turbine, high pressure & 511 & 1.9 \\
HTHP regenerator R1 & 320 & 1.2 \\
ORC pumps, mixer, flash separator & 12 & $<0.1$ \\
\midrule
$\sum I_j$ & 26415 & 100.0 \\
\bottomrule
\end{tabular}
\end{table}

First, heat transfer dominates. The five heat exchangers that carry the
process duty (HTHP condenser and evaporator, ORC evaporator and condenser, and
the borehole field) account for \SI{48.0}{\percent} of the destruction, the
compressors and turbines for \SI{18.5}{\percent}, and the motor and generator
for \SI{8.2}{\percent}. Second, the throttle between the condenser and the flash
separator destroys more exergy than either compressor stage, which indicates
that an expander in place of that valve is worth more than any plausible gain
in compressor efficiency. Third, the borehole field is responsible for only
\SI{6.8}{\percent} of the destruction; the case for the cascade and for tight
approach temperatures rests on the well count and the hardware rather than on
lost work \chk{$I_\mathrm{BB}$ is not yet independently verified}.

Pricing the geothermal stream lowers the headline efficiency from RTE $=0.317$
to $\psi=0.236$; under the stripped convention $\psi_\mathrm{str}=0.276$. The
geothermal flow required is \SI{100.8}{\kilo\gram\per\second}, i.e.\ 8.1
producers at the provisional yield, against 15.6 storage wells, and
\SI{57}{\percent} of the exergy lifted is reinjected.

\paragraph{Cooling of the source} $\Delta T_\mathrm{3A,4A}$ sets both the HTHP
evaporating temperature, and hence $\COP$, and the geothermal flow, which falls
much faster. Table~\ref{tab:dt3a4a} shows the consequence: the stripped
efficiency decreases monotonically, favouring the smallest cooling and the
largest geothermal flow, whereas the resource efficiency has an interior
maximum near \SI{20}{\kelvin}. At \SI{20}{\kelvin} the COP falls by
\SI{7.8}{\percent} relative to the reference case, while the geothermal flow
falls to \SI{47}{\percent} of its value: $N_\mathrm{geo}$ drops from
\num{8.07} to \num{3.80} and the whole field from \num{22.7} to
\SI{19.4}{\per\mega\watt}, for an RTE of \num{0.293}. The reference
value of \SI{10}{\kelvin} is therefore poorly placed, and the upper limit on
$\Delta T_\mathrm{3A,4A}$ is set by the formation constraint
$T_\mathrm{reinj,min}$, for which no measured value is yet available.

\begin{table}[htbp]
\centering
\caption{Effect of the cooling of the geothermal source. Every row is a
full cyclic-steady-state run on the basis of Eq.~\eqref{eq:psi}, with the
parasitics charged, so the rows are directly comparable with
Table~\ref{tab:ref-results}.}
\label{tab:dt3a4a}
\small
\begin{tabular}{@{}r r r r r r r@{}}
\toprule
$\Delta T_\mathrm{3A,4A}$ [K] & $\COP$ & RTE & $\md_\mathrm{geo}/\md_\mathrm{geo}^{(5\,\mathrm{K})}$ &
$u_\mathrm{geo}$ & $\psi_\mathrm{str}$ & $\psi$ \\
\midrule
5  & 2.582 & 0.3310 & 1.000 & 0.227 & 0.2833 & 0.1899 \\
10 & 2.473 & 0.3172 & 0.486 & 0.426 & 0.2764 & 0.2356 \\
15 & 2.373 & 0.3045 & 0.315 & 0.597 & 0.2698 & 0.2505 \\
20 & 2.280 & 0.2927 & 0.229 & 0.739 & 0.2634 & \textbf{0.2544} \\
25 & 2.194 & 0.2818 & 0.178 & 0.852 & 0.2571 & 0.2533 \\
30 & 2.114 & 0.2716 & 0.144 & 0.933 & 0.2511 & 0.2498 \\
\bottomrule
\end{tabular}
\end{table}

\subsection{Parametric assessment}
\label{sec:res-param}

\todo{Following the IHTC-18 paper: $\Nlay$ and the glide. Report RTE,
$\Nw$ (or $\mathcal D$ at fixed $\Nw$), $\bar\varepsilon_\mathrm{dc}$ and total
$UA$ on common axes. Expected interior optimum in $\Delta T_\mathrm{3D,2D}$
from the pinch, Eq.~\eqref{eq:pinch-necessary}, and in
$\Delta T_\mathrm{casc}$ subject to Eq.~\eqref{eq:span-bound}.}
\todo{$\mathcal D$ and $\bar\varepsilon_\mathrm{dc}$ versus $\Nw$ to separate
rate from inventory limitation.}
\todo{Carry $\psi$, $N_\mathrm{geo}$ and $(\Nw+N_\mathrm{geo})/\dot
W_\mathrm{el,out}$ in every parametric table; a study ranked on RTE alone
selects the free-source corner.}

\subsection{Limitations}
\label{sec:limitations}

\todo{Quantify, at least by estimate:}
(i)~H3 is at its limit at the reference case: the fronts of adjacent legs
meet over most of the well, so the residual solid in the corners between legs
is reached through a longer and narrower path than the annulus represents,
and the PCM-side resistance is understated;
(ii)~H9 neglects conduction between the two legs of a hairpin, which differ
by up to the cascade span at the wellhead;
(iii)~H6 neglects heat exchange with the formation, so the storage efficiency
is unity by construction and the store appears in the exergy audit only through
its internal destruction;
(iv)~the cycles have no pressure drop or heat loss in the exchangers and a
zero-approach ORC feed heater, and the isentropic efficiencies are constant;
(v)~the model has not been validated against experiment or against a
higher-fidelity (two-dimensional) conduction calculation. Items (i) and (ii)
are geometric questions about the cell boundary and would be settled together
by a two-dimensional conduction model of the real cross-section.

%=======================================================================
\section{Conclusions}
\label{sec:conclusions}
%=======================================================================
\todo{After Section~\ref{sec:results}.}

\section*{Acknowledgements}
\todo{Funding.}

%=======================================================================
\section*{Nomenclature}
%=======================================================================
\begingroup\small
\begin{longtable}{@{}p{0.16\linewidth}p{0.62\linewidth}p{0.16\linewidth}@{}}
$A$ & area; PCM cross-section per unit length & \si{\meter\squared}\\
$\Aav$ & PCM cross-section per hairpin per unit developed length & \si{\meter\squared}\\
$\Amelt$ & melted cross-section & \si{\meter\squared}\\
$C_s$, $C_l$ & sensible capacity per unit length & \si{\joule\per\meter\per\kelvin}\\
$c_p$ & specific heat & \si{\joule\per\kilogram\per\kelvin}\\
$D$ & diameter & \si{\meter}\\
$\mathcal D$ & deviation of delivered energy from target & --\\
$E'$ & PCM enthalpy per unit length & \si{\joule\per\meter}\\
$\Elat$, $\Eeq$ & latent capacity; equilibrium enthalpy per unit length & \si{\joule\per\meter}\\
$\Delta E$ & energy per cycle & \si{\joule}, \si{\kilo\watt\hour}\\
$f$ & Darcy friction factor & --\\
$h$ & specific enthalpy & \si{\joule\per\kilogram}\\
$h_i$, $h_e$ & internal, PCM-shell heat transfer coefficient & \si{\watt\per\meter\squared\per\kelvin}\\
$\hm$ & latent heat of fusion & \si{\joule\per\kilogram}\\
$K$ & segment conductance per unit length & \si{\watt\per\meter\per\kelvin}\\
$k$ & thermal conductivity & \si{\watt\per\meter\per\kelvin}\\
$L$ & length & \si{\meter}\\
$\md$ & mass flow rate per hairpin & \si{\kilogram\per\second}\\
$n_f$, $n_s$, $n_t$ & fins per leg, segments, hairpins per well & --\\
$\Nlay$, $\Nw$ & PCM layers, wells & --\\
$P_T$ & finned outer perimeter & \si{\meter}\\
$p$ & pressure & \si{\pascal}\\
$\dot Q$, $q'$ & heat rate; heat rate per unit length & \si{\watt}, \si{\watt\per\meter}\\
$R''_{f,i}$ & internal fouling resistance & \si{\meter\squared\kelvin\per\watt}\\
$r_i$, $r_e$, $r_\mathrm{cell}$ & tube inner, outer, and cell radius & \si{\meter}\\
$S$ & mean-to-surface shape factor & --\\
$s$ & developed coordinate; specific entropy & \si{\meter}; \si{\joule\per\kilogram\per\kelvin}\\
$T$ & temperature & \si{\kelvin}\\
$t$ & time; fin thickness & \si{\second}; \si{\meter}\\
$\Ui$ & overall coefficient on the inner area & \si{\watt\per\meter\squared\per\kelvin}\\
$UA$ & exchanger conductance & \si{\watt\per\kelvin}\\
$V$ & volume & \si{\meter\cubed}\\
$\dot W$, $W$ & power; work or electrical energy per cycle & \si{\watt}; \si{\kilo\watt\hour}\\
$\Ex$, $\Exd$ & exergy per cycle; exergy rate & \si{\kilo\watt\hour}; \si{\watt}\\
$I$ & exergy destruction per cycle & \si{\kilo\watt\hour}\\
$N_\mathrm{geo}$ & geothermal producing wells & --\\
$s_\mathrm{gen}$ & specific entropy generation & \si{\joule\per\kilogram\per\kelvin}\\
$T_0$ & dead-state temperature, $T_\mathrm{sink}$ & \si{\kelvin}\\
$u_\mathrm{geo}$ & fraction of the geothermal exergy used & --\\
$x_\mathrm{8h}$ & vapour fraction leaving the flash valve & --\\
$y$ & ORC reheat mass fraction & --\\[3pt]
\multicolumn{3}{@{}l}{\emph{Greek}}\\
$\beta$ & $r_\mathrm{cell}/r_e$ & --\\
$\Delta s$ & segment length & \si{\meter}\\
$\delta$ & shell thickness & \si{\meter}\\
$\varepsilon$ & melted fraction; effectiveness & --\\
$\eta$ & efficiency; $\eta_{s,t}$, $\eta_{s,p}$, $\eta_{s,c}$ isentropic efficiency of turbines, pumps, compressors & --\\
$\lambda$ & storage-loss allowance & --\\
$\rho$, $\rho_E$ & density; thermal energy density & \si{\kilogram\per\meter\cubed}; \si{\kilo\watt\hour\per\meter\cubed}\\
$\tau$ & relaxation time $C/K$ & \si{\second}\\
$\phi$ & fraction of ORC heat absorbed isothermally & --\\
$\psi$, $\psi_\mathrm{str}$ & exergetic efficiency, resource and stripped conventions & --\\[3pt]
\multicolumn{3}{@{}l}{\emph{Subscripts}}\\
\multicolumn{3}{@{}p{0.95\linewidth}@{}}{a source stream;
c, d secondary water during charging, discharging; ch, dc charging,
discharging; e ORC working fluid; geo geothermal; h HTHP refrigerant; l liquid;
m melting; reinj reinjected; s solid; w water; wf working fluid.}\\
\end{longtable}
\endgroup
